% ===========================================================
%  Глава 13 — GRPO и онлайн-RL с проверяемыми вознаграждениями
% ===========================================================
\chapter[GRPO и онлайн-RL]{GRPO и онлайн-RL с проверяемыми вознаграждениями}
\label{ch:grpo}

\begin{quote}
\itshape
Лучший способ проверить правильность ответа\,---\,проверить его, а не спрашивать кого-то, выглядит ли он правильным.
\end{quote}

\bigskip

В главе~\ref{ch:actor-critic} была разработана проксимальная оптимизация стратегии (PPO)\,---\,основной алгоритм обучения с подкреплением на основе обратной связи от человека (RLHF). PPO обеспечивает стабильные обновления стратегии с помощью отсечённой суррогатной цели и вычисляет преимущества через обученного критика (функцию ценности). В главе~\ref{ch:dpo} затем было показано, что прямая оптимизация предпочтений (DPO) может полностью обойтись без модели вознаграждения, параметризуя KL-регуляризованную цель в виде контролируемой функции потерь на парах предпочтений. DPO прост и стабилен, но принципиально автономен: он обучается на фиксированном наборе данных и не может исследовать пространство.

В данной главе вводится \textbf{групповая относительная оптимизация стратегии} (GRPO)\index{GRPO}\index{Group Relative Policy Optimisation|see{GRPO}}\,---\,алгоритм, лежащий в основе DeepSeek-R1 и других недавних моделей рассуждений (Shao и др., 2024; DeepSeek-AI, 2025). GRPO занимает удачную нишу между PPO и DPO: он является \emph{онлайн}\,---\,стратегия генерирует свежие завершения на каждой итерации,\,---\,но устраняет критика с функцией ценности, делающего PPO дорогостоящим. Вместо этого GRPO оценивает преимущества, сэмплируя \emph{группу} завершений для каждого промпта и нормируя вознаграждения внутри группы. Результатом является простой, экономичный по памяти алгоритм, сохраняющий преимущества стабильности механизма отсечения PPO.

Мы также формализуем \textbf{обучение с подкреплением с проверяемыми вознаграждениями} (RLVR)\index{RLVR}\index{Reinforcement Learning with Verifiable Rewards|see{RLVR}}\,---\,парадигму обучения, при которой сигнал вознаграждения поступает от детерминированного верификатора (например, набора юнит-тестов или символического математического проверщика), а не от обученной модели вознаграждения. RLVR является естественным дополнением GRPO: когда вознаграждения бинарны и не зашумлены, групп-относительное преимущество ведёт себя особенно хорошо, а отсутствие обученной модели вознаграждения устраняет риск взлома вознаграждения.

В разделе~\ref{sec:grpo-motivation} мотивируется GRPO путём выявления конкретных ограничений PPO и DPO, которые он устраняет. В разделе~\ref{sec:grpo-objective} выводится цель GRPO в полном объёме. В разделе~\ref{sec:grpo-algorithm} представлен полный алгоритм. В разделе~\ref{sec:grpo-theory} даётся теоретический анализ, включая связь с REINFORCE Leave-One-Out (RLOO, раздел~\ref{sec:rloo} главы~\ref{ch:pg}) и анализ смещения-дисперсии размера группы. В разделе~\ref{sec:rlvr} формализуется RLVR. В разделе~\ref{sec:rejection-sampling} обсуждается отбор по отклонению и best-of-$N$ как более простой базовый уровень. В разделе~\ref{sec:grpo-scale} описывается инженерия онлайн-обучения RL в масштабе. В разделе~\ref{sec:grpo-comparison} сравниваются PPO, GRPO и DPO в единой таблице. В разделе~\ref{sec:grpo-summary} подводятся итоги главы.

% ===========================================================
\section{Мотивация: за пределами PPO и DPO}
\label{sec:grpo-motivation}

\subsection{Стоимость критика в PPO}

\index{PPO!critic cost}
Напомним из раздела~\ref{sec:ppo}, что отсечённая суррогатная цель PPO~\eqref{eq:ppo-clip} требует оценок преимущества $\hat{A}_t$, вычисляемых через обобщённую оценку преимущества (GAE, определение~\ref{def:gae}). GAE, в свою очередь, требует обученной функции ценности $V_\phi(S_t)$\,---\,\emph{критика},\,---\,параметры $\phi$ которого обучаются вместе с актором. В рамках RLHF, описанного в разделе~\ref{sec:ppo-llm}, архитектура из четырёх моделей (актор, критик, референсная модель, модель вознаграждения) требует огромного объёма видеопамяти GPU. При масштабе 70B параметров каждая копия модели занимает примерно 140 ГБ в формате половинной точности, так что четыре копии требуют 560 ГБ\,---\,больше, чем на одном узле 8$\times$A100.

Критик также является источником нестабильности. Ему приходится отслеживать движущуюся цель: по мере улучшения актора $V^{\pi_\theta}$ меняется, и критик должен переобучаться. В начале обучения критик особенно неточен, производя зашумлённые оценки преимуществ, которые могут дестабилизировать актора. Тщательная инициализация (раздел~\ref{sec:ppo-llm}) смягчает, но не устраняет эту проблему.

\subsection{Проблема сдвига распределения в DPO}

\index{DPO!distribution shift}
В главе~\ref{ch:dpo} было показано, что DPO устраняет как модель вознаграждения, так и критика, сокращая требования к памяти до двух копий моделей ($\pi_\theta$ и $\pi_\mathrm{ref}$). Платой является то, что DPO является автономным: он обучается на фиксированном наборе данных предпочтений $\mathcal{D} = \{(x^{(i)}, y_w^{(i)}, y_l^{(i)})\}$, собранном при стратегии SFT. По мере обучения и удаления $\pi_\theta$ от $\pi_\mathrm{SFT}$ метки предпочтений устаревают. Ответы, которые улучшенная стратегия теперь генерировала бы, отсутствуют в $\mathcal{D}$, и модель не может их открыть. В разделе~\ref{sec:dpo-limitations} главы~\ref{ch:dpo} это ограничение было проанализировано подробно: DPO сходится к оптимальной стратегии \emph{для данного фиксированного набора данных}, но дальнейшее улучшение требует новых данных.

\subsection{GRPO: лучшее из двух миров}

\index{GRPO!motivation}
GRPO одновременно решает обе проблемы:
\begin{description}[leftmargin=!,labelwidth=3.5cm]
  \item[Без критика] Преимущества оцениваются по группе сэмплированных завершений, а не по обученной функции ценности. Объём занимаемой памяти сокращается с четырёх до двух копий модели (актор $\pi_\theta$ и референс $\pi_\mathrm{ref}$), как и в DPO.
  \item[Онлайн-режим] Стратегия генерирует новые завершения на каждой итерации и получает свежую обратную связь по вознаграждению. Сдвига распределения нет: сигнал вознаграждения всегда откалиброван под выходное распределение текущей стратегии.
  \item[Стабильные обновления] GRPO унаследует отсечённое отношение важности PPO, предотвращая чрезмерно крупные шаги обновления стратегии.
\end{description}
Ключевое понимание состоит в том, что основная роль критика\,---\,обеспечение базиса для снижения дисперсии\,---\,может быть выполнена эмпирическими статистиками группы завершений, сэмплированных из текущей стратегии. Эта идея тесно связана с оценщиком RLOO (раздел~\ref{sec:rloo} главы~\ref{ch:pg}), который использует базис leave-one-out; GRPO расширяет его стандартизацией и отсечением в стиле PPO.

На рисунке~\ref{fig:grpo-architecture} сравниваются архитектуры модели PPO-RLHF (четыре модели) и GRPO-RLVR (две модели плюс верификатор).

\begin{figure}[ht]
\centering
\begin{tikzpicture}[
    node distance=1.2cm and 2.0cm,
    box/.style={rectangle, draw, rounded corners=4pt, minimum width=2.2cm,
                minimum height=0.8cm, font=\small, fill=blue!8},
    verifier/.style={rectangle, draw, rounded corners=4pt, minimum width=2.2cm,
                minimum height=0.8cm, font=\small, fill=green!12},
    env/.style={rectangle, draw, rounded corners=4pt, minimum width=2.2cm,
                minimum height=0.8cm, font=\small, fill=orange!12},
    arr/.style={-Stealth, thick},
    lab/.style={font=\footnotesize, inner sep=2pt}
  ]
  % Сторона PPO
  \node[font=\small\bfseries] at (0, 3.2) {PPO-RLHF (4 модели)};
  \node[box] (actor1) at (-1.2, 2.2) {Актор $\pi_\theta$};
  \node[box] (critic) at (1.2, 2.2) {Критик $V_\phi$};
  \node[box] (ref1) at (-1.2, 0.8) {Реф.\ $\pi_\mathrm{ref}$};
  \node[env] (rm) at (1.2, 0.8) {Модель вознаграждения $r_\psi$};

  % Сторона GRPO
  \node[font=\small\bfseries] at (5.5, 3.2) {GRPO-RLVR (2 модели)};
  \node[box] (actor2) at (4.3, 2.2) {Актор $\pi_\theta$};
  \node[box] (ref2) at (6.7, 2.2) {Реф.\ $\pi_\mathrm{ref}$};
  \node[verifier] (ver) at (5.5, 0.8) {Верификатор $V(q,o)$};

  % Метки
  \node[lab, below=0.15cm of critic, gray] {(исключено)};
  \node[lab, below=0.15cm of rm, gray] {(заменено верификатором)};
\end{tikzpicture}
\caption{Сравнение архитектур моделей. PPO-RLHF требует четырёх копий нейронных сетей: актора, критика, референса и модели вознаграждения. GRPO-RLVR требует только двух копий (актора и референса) плюс лёгкий детерминированный верификатор, уменьшая вдвое объём занимаемой памяти.}
\label{fig:grpo-architecture}
\end{figure}

% ===========================================================
\section{Цель GRPO}
\label{sec:grpo-objective}

\index{GRPO!history}
GRPO был предложен Shao и др. (2024) в статье DeepSeekMath, где его использовали для улучшения математических рассуждений в языковой модели с 7B параметрами. Алгоритм был затем принят в качестве компонента RL в конвейере DeepSeek-R1 (DeepSeek-AI, 2025), достигнув результатов уровня лучших мировых решений по математике уровня соревнований, генерации кода и решению задач по естественным наукам. Название «групп-относительный» отражает две определяющие особенности: преимущества вычисляются относительно \emph{группы} завершений для одного и того же промпта и выражаются как \emph{относительные} отклонения (стандартизованные групповым стандартным отклонением).

\subsection{Постановка задачи и обозначения}

\index{GRPO!notation}
Пусть $\pi_\theta$ обозначает обучаемую стратегию языковой модели, $\pi_{\theta_\mathrm{old}}$\,---\,стратегию в начале текущей итерации (до шага градиента), а $\pi_\mathrm{ref}$\,---\,замороженную референсную стратегию (обычно контрольную точку SFT). Промпт $q$ берётся из набора данных $\mathcal{D}$. Для каждого промпта GRPO сэмплирует \textbf{группу}\index{group sampling} из $G$ независимых завершений из старой стратегии:
\begin{equation}
\label{eq:grpo-sampling}
o_1, o_2, \ldots, o_G \;\sim\; \pi_{\theta_\mathrm{old}}(\cdot \mid q),
\end{equation}
где каждое завершение $o_i = (o_{i,1}, o_{i,2}, \ldots, o_{i,|o_i|})$ является последовательностью токенов переменной длины. Функция вознаграждения $r(q, o_i) \in \R$ оценивает каждое завершение, давая вознаграждения $r_1, r_2, \ldots, r_G$.

\subsection{Групп-относительное преимущество}

\index{GRPO!advantage}
В PPO преимущество $\hat{A}_t$ вычисляется через обученного критика с помощью GAE. GRPO заменяет это \textbf{групп-относительным преимуществом}\index{group-relative advantage}, не требующим обученных параметров:

\begin{definition}[Преимущество GRPO]
\label{def:grpo-advantage}
Для завершения $i$ в группе из $G$ завершений с вознаграждениями $r_1, \ldots, r_G$ преимущество GRPO равно
\begin{equation}
\label{eq:grpo-advantage}
\hat{A}_i = \frac{r_i - \mu_G}{\sigma_G + \varepsilon_0},
\end{equation}
где $\mu_G = \frac{1}{G}\sum_{j=1}^G r_j$\,---\,среднее по группе, $\sigma_G = \sqrt{\frac{1}{G}\sum_{j=1}^G (r_j - \mu_G)^2}$\,---\,стандартное отклонение по группе, а $\varepsilon_0 > 0$\,---\,малая константа для численной стабильности (обычно $10^{-8}$).
\end{definition}

Преимущество $\hat{A}_i$ измеряет, на сколько стандартных отклонений выше или ниже среднего по группе находится вознаграждение завершения $i$. Если все завершения в группе получают одинаковое вознаграждение (например, все правильные или все неправильные), то $\sigma_G = 0$ и $\hat{A}_i \approx 0$ для всех $i$, не производя никакого градиентного сигнала. Это желательно: промпт, который модель всегда решает или никогда не решает, не даёт никакой контрастной информации.

\begin{remark}
Одно и то же преимущество $\hat{A}_i$ присваивается \emph{каждому токену} в завершении $i$:
\[
\hat{A}_{i,t} = \hat{A}_i \quad \mbox{для всех } t \in \{1, \ldots, |o_i|\}.
\]
Это преимущество на уровне последовательности (или уровне результата), а не на уровне токена. В отличие от GAE, который присваивает разные преимущества разным токенам на основе оценок функции ценности на каждом шаге, GRPO равномерно приписывает итоговое вознаграждение всем токенам, которые его произвели. Это упрощение является одновременно достоинством (критик не нужен) и ограничением (нет присваивания заслуг внутри завершения).
\end{remark}

\subsection{Отсечённая суррогатная цель}

\index{GRPO!clipped objective}
Следуя PPO (определение~\ref{def:ppo-clip}), GRPO использует взвешенную по важности цель с отсечением. Для завершения $i$ и позиции токена $t$ определим отношение важности на уровне токена:
\begin{equation}
\label{eq:grpo-ratio}
\rho_{i,t}(\theta) = \frac{\pi_\theta(o_{i,t} \mid q, o_{i,<t})}{\pi_{\theta_\mathrm{old}}(o_{i,t} \mid q, o_{i,<t})}.
\end{equation}
Отсечённый суррогат для данного токена равен
\begin{equation}
\label{eq:grpo-clip-token}
L_{i,t}^{\mathrm{CLIP}}(\theta) = \min\!\Big(\rho_{i,t}(\theta)\,\hat{A}_i,\; \mathrm{clip}\!\big(\rho_{i,t}(\theta),\, 1-\epsilon,\, 1+\epsilon\big)\,\hat{A}_i\Big),
\end{equation}
где $\epsilon > 0$\,---\,диапазон отсечения (обычно $\epsilon \in [0{,}1, 0{,}3]$). Механизм отсечения идентичен PPO: он ограничивает стимул увеличивать вероятность положительно-адвантажных завершений и ограничивает стимул уменьшать вероятность отрицательно-адвантажных завершений, предотвращая дестабилизирующие крупные шаги стратегии.

\subsection{Штраф KL}

\index{GRPO!KL penalty}
GRPO добавляет штраф дивергенции Кульбака\,---\,Лейблера непосредственно в функцию потерь, а не встраивает его в вознаграждение, как это делает PPO (уравнение~\eqref{eq:rlhf-reward}). Такая конструкция отделяет KL-регуляризацию от вычисления преимущества, избегая осложнений при групповой нормализации. Оценщик KL на уровне токена, используемый на практике, равен
\begin{equation}
\label{eq:grpo-kl}
D_{\mathrm{KL}}^{(i,t)} = \frac{\pi_\mathrm{ref}(o_{i,t} \mid q, o_{i,<t})}{\pi_\theta(o_{i,t} \mid q, o_{i,<t})} - \log\frac{\pi_\mathrm{ref}(o_{i,t} \mid q, o_{i,<t})}{\pi_\theta(o_{i,t} \mid q, o_{i,<t})} - 1,
\end{equation}
что является несмещённым, неотрицательным оценщиком Шульмана (2020) дивергенции KL между $\pi_\theta$ и $\pi_\mathrm{ref}$. Этот оценщик всегда удовлетворяет $D_{\mathrm{KL}}^{(i,t)} \geq 0$, равенство достигается тогда и только тогда, когда $\pi_\theta(o_{i,t} \mid q, o_{i,<t}) = \pi_\mathrm{ref}(o_{i,t} \mid q, o_{i,<t})$.

\begin{remark}
Альтернативным оценщиком является более простой $\log[\pi_\theta(o_{i,t} \mid q, o_{i,<t}) / \pi_\mathrm{ref}(o_{i,t} \mid q, o_{i,<t})]$, который несмещён, но может быть отрицательным для отдельных сэмплов. Оценщик в~\eqref{eq:grpo-kl} предпочтительнее, поскольку его неотрицательность стабилизирует обучение.
\end{remark}

\subsection{Полная цель GRPO}

\index{GRPO!objective}
Объединяя отсечённый суррогат и штраф KL, цель GRPO принимает вид:

\begin{definition}[Цель GRPO]
\label{def:grpo-objective}
\begin{equation}
\label{eq:grpo-objective}
\boxed{
J_{\mathrm{GRPO}}(\theta)
= \E_{\substack{q \sim \mathcal{D} \\ \{o_i\}_{i=1}^G \sim \pi_{\theta_\mathrm{old}}(\cdot|q)}}
\!\left[
\frac{1}{G}\sum_{i=1}^{G} \frac{1}{|o_i|}\sum_{t=1}^{|o_i|}
\Big(
  L_{i,t}^{\mathrm{CLIP}}(\theta)
  - \beta\, D_{\mathrm{KL}}^{(i,t)}
\Big)
\right],
}
\end{equation}
где $L_{i,t}^{\mathrm{CLIP}}$ определяется уравнением~\eqref{eq:grpo-clip-token}, $D_{\mathrm{KL}}^{(i,t)}$\,---\,уравнением~\eqref{eq:grpo-kl}, а $\hat{A}_i$\,---\,уравнением~\eqref{eq:grpo-advantage}.
\end{definition}

Ряд особенностей этой цели заслуживает комментария.

\begin{itemize}[leftmargin=*]
\item Внутренняя сумма усредняет по позициям токенов внутри каждого завершения. Нормировка на $1/|o_i|$ предотвращает непропорционально большой вклад длинных завершений в градиент.
\item Внешняя сумма усредняет по $G$ завершениям в группе. Это аналогично усреднению по $K$ траекториям в стандартном мини-батче градиента стратегии, но здесь траектории разделяют один и тот же промпт.
\item Коэффициент KL $\beta > 0$ управляет дрейфом от референсной стратегии. Большее $\beta$ сильнее ограничивает стратегию, снижая риск взлома вознаграждения ценой более медленного обучения.
\item Когда на одном и том же батче завершений делается несколько шагов градиента, $\pi_{\theta_\mathrm{old}}$ фиксируется при стратегии, сгенерировавшей завершения, и обновляется только $\pi_\theta$. Это та же коррекция вне текущей стратегии, что и в PPO.
\end{itemize}

\begin{proposition}[Локальная эквивалентность градиентов]
\label{prop:grpo-local-gradient}
При $\theta = \theta_\mathrm{old}$ градиент $J_{\mathrm{GRPO}}(\theta)$ по $\theta$ равен градиенту стратегии REINFORCE с базисом среднего по группе:
\begin{equation}
\label{eq:grpo-local-grad}
\nabla_\theta J_{\mathrm{GRPO}}(\theta)\big|_{\theta = \theta_\mathrm{old}}
= \frac{1}{G}\sum_{i=1}^G \hat{A}_i \sum_{t=1}^{|o_i|} \frac{1}{|o_i|}\nabla_\theta \log \pi_\theta(o_{i,t} \mid q, o_{i,<t})\Big|_{\theta_\mathrm{old}}
- \beta\,\nabla_\theta D_{\mathrm{KL}}.
\end{equation}
\end{proposition}

\begin{proof}
При $\theta = \theta_\mathrm{old}$ отношение важности $\rho_{i,t}(\theta) = 1$ для всех $i, t$. Поскольку $1 \in [1-\epsilon, 1+\epsilon]$, отсечение не активно, поэтому $L_{i,t}^{\mathrm{CLIP}} = \rho_{i,t}\hat{A}_i$. Дифференцирование $\rho_{i,t}\hat{A}_i$ по $\theta$ при $\rho_{i,t} = 1$ даёт $\hat{A}_i \nabla_\theta \log \pi_\theta(o_{i,t} \mid q, o_{i,<t})$\,---\,в точности градиент REINFORCE, взвешенный групп-нормированным преимуществом.
\end{proof}

% ===========================================================
\section{Алгоритм GRPO}
\label{sec:grpo-algorithm}

Алгоритм~\ref{alg:grpo} резюмирует цикл обучения GRPO.

\begin{algorithm}[ht]
\caption{Групповая относительная оптимизация стратегии (GRPO)}
\label{alg:grpo}
\DontPrintSemicolon
\KwIn{Стратегия $\pi_\theta$, референсная стратегия $\pi_\mathrm{ref}$, функция вознаграждения $r$, набор данных промптов $\mathcal{D}$}
\KwIn{Размер группы $G$, диапазон отсечения $\epsilon$, коэффициент KL $\beta$, скорость обучения $\alpha$, число внутренних эпох $K_e$}
\KwOut{Обученная стратегия $\pi_\theta$}
Инициализировать $\pi_\theta \leftarrow \pi_\mathrm{ref}$\;
\Repeat{сходимость}{
  Сэмплировать батч промптов $\{q^{(1)}, \ldots, q^{(B)}\}$ из $\mathcal{D}$\;
  \tcp{Фаза генерации}
  \For{каждого промпта $q^{(b)}$, $b = 1, \ldots, B$}{
    Сэмплировать $G$ завершений $o_1^{(b)}, \ldots, o_G^{(b)} \sim \pi_\theta(\cdot \mid q^{(b)})$\;
    Вычислить вознаграждения $r_i^{(b)} \leftarrow r(q^{(b)}, o_i^{(b)})$ для $i = 1, \ldots, G$\;
    Вычислить групповое среднее $\mu_G^{(b)} \leftarrow \frac{1}{G}\sum_{i=1}^G r_i^{(b)}$ и std $\sigma_G^{(b)}$\;
    Вычислить преимущества $\hat{A}_i^{(b)} \leftarrow (r_i^{(b)} - \mu_G^{(b)}) / (\sigma_G^{(b)} + \varepsilon_0)$ для $i = 1, \ldots, G$\;
    Сохранить логарифмические вероятности $\log \pi_{\theta_\mathrm{old}}(o_{i,t}^{(b)} \mid q^{(b)}, o_{i,<t}^{(b)})$ для всех $i, t$\;
  }
  \tcp{Фаза обучения}
  Зафиксировать $\theta_\mathrm{old} \leftarrow \theta$\;
  \For{эпоха $= 1, \ldots, K_e$}{
    \For{каждого мини-батча пар промпт-завершение}{
      Вычислить $\rho_{i,t} \leftarrow \pi_\theta(o_{i,t} \mid q, o_{i,<t})\, /\, \pi_{\theta_\mathrm{old}}(o_{i,t} \mid q, o_{i,<t})$ для всех $i, t$\;
      Вычислить отсечённые цели $L_{i,t}^{\mathrm{CLIP}}$ по~\eqref{eq:grpo-clip-token}\;
      Вычислить штрафы KL $D_{\mathrm{KL}}^{(i,t)}$ по~\eqref{eq:grpo-kl}\;
      Вычислить потерю $\mathcal{L} \leftarrow -\frac{1}{|\mbox{батч}|}\sum_{i}\frac{1}{|o_i|}\sum_t (L_{i,t}^{\mathrm{CLIP}} - \beta\, D_{\mathrm{KL}}^{(i,t)})$\;
      $\theta \leftarrow \theta - \alpha\,\nabla_\theta \mathcal{L}$\;
    }
  }
  \tcp{Опционально обновить референсную модель}
  Каждые $N_\mathrm{ref}$ шагов: $\pi_\mathrm{ref} \leftarrow \pi_\theta$\;
}
\Return $\pi_\theta$
\end{algorithm}

\subsection{Ключевые гиперпараметры}

\begin{description}[leftmargin=!,labelwidth=3.5cm]
  \item[Размер группы $G$] Число завершений на промпт. Типичные значения: $G \in \{4, 8, 16, 64\}$. Большее $G$ обеспечивает оценки преимущества с меньшей дисперсией, но линейно увеличивает стоимость генерации. DeepSeek-R1 использует $G = 64$ в начале обучения.
  \item[Диапазон отсечения $\epsilon$] Управляет доверительной областью. Стандартный выбор $\epsilon = 0{,}2$, как и в PPO.
  \item[Коэффициент KL $\beta$] Управляет дрейфом от референса. Обычные значения от $\beta = 0{,}01$ (свободный) до $\beta = 0{,}1$ (жёсткий). Некоторые реализации используют динамическое $\beta$ с целевым бюджетом KL.
  \item[Число внутренних эпох $K_e$] Число шагов градиента на батч сгенерированных завершений. DeepSeek-R1 использует $K_e = 1$ (один проход по каждому батчу)\,---\,более консервативный подход, чем типичное $K_e \in [3, 10]$ в PPO.
  \item[Частота обновления референса $N_\mathrm{ref}$] Как часто референсная модель заменяется текущей стратегией. DeepSeek-R1 обновляет каждые 400 шагов. Частые обновления позволяют штрафу KL адаптироваться к улучшенной стратегии; редкие обновления обеспечивают более стабильный якорь.
\end{description}

\subsection{Проработанный пример}

\begin{example}[Вычисление преимущества GRPO]
\label{ex:grpo-advantage}
Рассмотрим промпт $q$ = «Чему равно $\int_0^1 x^2\,dx$?» с $G = 4$ завершениями:
\begin{center}
{\footnotesize
\begin{tabular}{@{}clcc@{}}
\toprule
$i$ & Завершение & $r_i$ & $\hat{A}_i$ \\
\midrule
1 & «$\frac{1}{3}$. Интегрируем $x^2$, получаем $\frac{x^3}{3}$\ldots» (верно) & 1,0 & $+0{,}94$ \\
2 & «$\frac{1}{2}$. Интеграл равен $\frac{x^2}{2}$\ldots» (неверно) & 0,0 & $-1{,}09$ \\
3 & «$1/3$. По степенному правилу\ldots» (верно) & 1,0 & $+0{,}94$ \\
4 & «Не знаю.» (неверно) & 0,0 & $-1{,}09$ \\
\bottomrule
\end{tabular}
}
\end{center}
Статистики группы: $\mu_G = 0{,}5$ и $\sigma_G \approx 0{,}577$. По~\eqref{eq:grpo-advantage}: $\hat{A}_1 = (1{,}0 - 0{,}5)/0{,}577 \approx +0{,}87$; $\hat{A}_2 = (0{,}0 - 0{,}5)/0{,}577 \approx -0{,}87$. (Значения в таблице используют выборочное стандартное отклонение с поправкой Бесселя для иллюстрации.) Положительные преимущества подкрепляют правильные завершения, а отрицательные подавляют неправильные. Критик не потребовался.\qed
\end{example}

% ===========================================================
\section{Теоретический анализ}
\label{sec:grpo-theory}

\subsection{Связь с REINFORCE с базисом}

\index{GRPO!connection to REINFORCE}
В утверждении~\ref{prop:grpo-local-gradient} было показано, что при $\theta = \theta_\mathrm{old}$ градиент GRPO сводится к REINFORCE со средним по группе в качестве базиса. Для точности рассмотрим «неотсечённую GRPO» (при $\epsilon \to \infty$):
\begin{equation}
\label{eq:grpo-unclipped}
J_{\mathrm{GRPO}}^{\mbox{без отсечения}}(\theta) = \frac{1}{G}\sum_{i=1}^G \hat{A}_i \cdot \frac{1}{|o_i|}\sum_{t=1}^{|o_i|} \log \pi_\theta(o_{i,t} \mid q, o_{i,<t}) - \beta\, \hat{D}_\mathrm{KL}.
\end{equation}
Игнорируя нормировку $1/|o_i|$ и член KL, первая часть равна
\[
\frac{1}{G}\sum_{i=1}^G \frac{r_i - \mu_G}{\sigma_G} \log \pi_\theta(o_i \mid q)
= \frac{1}{\sigma_G}\left[\frac{1}{G}\sum_{i=1}^G r_i \log \pi_\theta(o_i \mid q) - \mu_G \cdot \frac{1}{G}\sum_{i=1}^G \log \pi_\theta(o_i \mid q)\right].
\]
Первый член\,---\,стандартный оценщик REINFORCE. Второй член вычитает базис среднего по группе $\mu_G$, умноженный на среднюю логарифмическую вероятность. Множитель $1/\sigma_G$ масштабирует весь градиент, влияя на эффективную скорость обучения, но не на его направление.

\begin{theorem}[Несмещённость базиса среднего по группе]
\label{thm:grpo-unbiased}
\index{GRPO!unbiasedness}
Оценщик градиента GRPO (без отсечения и нормировки по длине) является несмещённым оценщиком градиента стратегии $\nabla_\theta J(\theta)$ с точностью до зависящего от промпта масштабного множителя $1/\sigma_G$.
\end{theorem}

\begin{proof}
Базис среднего по группе $\mu_G = \frac{1}{G}\sum_j r_j$ включает вознаграждение $r_i$ самого $i$-го сэмпла. В отличие от базиса RLOO (уравнение~\eqref{eq:rloo-baseline}), это делает $\mu_G$ \emph{слабо} зависимым от $o_i$. Однако вводимое смещение мало. Запишем $\mu_G = \frac{r_i}{G} + \frac{1}{G}\sum_{j \neq i} r_j$. Второй член $b_i' = \frac{1}{G}\sum_{j \neq i} r_j$ не зависит от $o_i$ при данном $q$ и, следовательно, является допустимым базисом по теореме о базисе (теорема~\ref{thm:baseline}). Дополнительный член $r_i / G$ вводит смещение порядка $O(1/G)$:
\[
\E\!\left[\frac{r_i}{G} \nabla_\theta \log \pi_\theta(o_i \mid q)\right] = \frac{1}{G}\nabla_\theta J(\theta),
\]
что является масштабированной версией истинного градиента. После вычитания из члена $r_i \nabla_\theta \log \pi_\theta$ чистый эффект состоит в том, что оценщик GRPO аппроксимирует $(1 - 1/G)\nabla_\theta J(\theta)$\,---\,смещённый на множитель $(G-1)/G$, что пренебрежимо мало при $G \geq 4$. Деление на $\sigma_G$ является промпт-зависимым скаляром и не влияет на направление градиента в среднем.
\end{proof}

\subsection{Сравнение с RLOO}

\index{GRPO!vs RLOO}
В замечании~\ref{rem:rloo-grpo} главы~\ref{ch:pg} отмечалась тесная связь между GRPO и RLOO. Таблица~\ref{tab:grpo-vs-rloo} делает это сравнение явным.

\begin{table}[ht]
\centering
\caption{Сравнение оценщиков преимущества RLOO и GRPO.}
\label{tab:grpo-vs-rloo}
\footnotesize
\begin{tabular}{@{}lll@{}}
\toprule
\textbf{Аспект} & \textbf{RLOO} & \textbf{GRPO} \\
\midrule
Базис & Среднее $K-1$ других сэмплов & Среднее всех $G$ сэмплов \\
Нормировка & Нет (сырое отклонение) & Деление на $\sigma_G$ \\
Смещение & Точно несмещён & Смещение $O(1/G)$ \\
Отсечение & Нет & Отсечение в стиле PPO \\
Размещение KL & В вознаграждении & В функции потерь \\
Много эпох & Не стандартно & Да (как PPO) \\
\bottomrule
\end{tabular}
\end{table}

Преимущество RLOO для сэмпла $i$ равно $\hat{A}_i^{\mathrm{RLOO}} = \frac{K}{K-1}(r_i - \bar{r})$ (уравнение~\eqref{eq:rloo-clean}), где $\bar{r} = \frac{1}{K}\sum_j r_j$. Преимущество GRPO равно $\hat{A}_i^{\mathrm{GRPO}} = (r_i - \bar{r}) / \sigma_G$. Ключевые различия: (i)~RLOO использует точный базис leave-one-out и точно несмещён, тогда как GRPO использует среднее по группе, включая сэмпл $i$, и вводит смещение $O(1/G)$; (ii)~GRPO делит на $\sigma_G$, что стандартизирует масштаб преимущества по промптам с разной дисперсией вознаграждения; (iii)~GRPO применяет отсечение в стиле PPO, обеспечивая безопасное обучение с несколькими эпохами на одном батче.

\subsection{Снижение дисперсии при групповой нормализации}

\index{GRPO!variance reduction}
\begin{proposition}[Дисперсия преимущества GRPO]
\label{prop:grpo-variance}
Пусть $r_1, \ldots, r_G$ независимы и одинаково распределены со средним $\mu$ и дисперсией $\sigma^2$. Дисперсия ненормированного преимущества $r_i - \mu_G$ равна
\begin{equation}
\label{eq:grpo-var-unnorm}
\Var[r_i - \mu_G] = \sigma^2\!\left(1 - \frac{1}{G}\right) = \frac{(G-1)\sigma^2}{G}.
\end{equation}
После стандартизации на $\sigma_G$ преимущество $\hat{A}_i$ имеет приблизительно единичную дисперсию независимо от распределения вознаграждений.
\end{proposition}

\begin{proof}
Поскольку $\mu_G = \frac{1}{G}\sum_j r_j$, имеем $r_i - \mu_G = r_i - \frac{1}{G}r_i - \frac{1}{G}\sum_{j \neq i} r_j = \frac{G-1}{G}r_i - \frac{1}{G}\sum_{j \neq i} r_j$. Используя независимость:
\[
\Var[r_i - \mu_G] = \left(\frac{G-1}{G}\right)^2 \sigma^2 + \frac{G-1}{G^2}\sigma^2 = \frac{(G-1)\sigma^2}{G}.
\]
Деление на $\sigma_G \approx \sigma\sqrt{(G-1)/G}$ (при больших $G$) даёт $\Var[\hat{A}_i] \approx 1$.
\end{proof}

Эта автоматическая нормализация дисперсии является практическим преимуществом GRPO перед RLOO: промпты с высокой дисперсией вознаграждения и промпты с низкой дисперсией вносят обновления градиента сопоставимого масштаба, не давая ни одному промпту доминировать в мини-батче.

\subsection{Влияние размера группы $G$}

\index{GRPO!group size}
Увеличение размера группы $G$ оказывает два эффекта:
\begin{enumerate}[leftmargin=*]
  \item \textbf{Лучшая оценка базиса.} Среднее по группе $\mu_G$ сходится к истинному ожидаемому вознаграждению $\E_{\pi_\theta}[r \mid q]$ со скоростью $O(1/\sqrt{G})$. Это снижает дисперсию оценки преимущества и улучшает качество градиента.
  \item \textbf{Более высокая стоимость генерации.} Каждый промпт требует $G$ прямых проходов через модель. Если модель генерирует в среднем $T$ токенов, общая стоимость генерации на промпт равна $O(GT)$, что может доминировать над стоимостью обучения.
\end{enumerate}
Оптимальное $G$ уравновешивает статистическую эффективность и вычислительные затраты. Эмпирически $G$ от 8 до 64 хорошо работает для задач рассуждения. Для бинарных вознаграждений (правильно / неправильно) более крупные группы особенно выгодны, поскольку они обеспечивают более разнообразные контрастные пары.

\begin{example}[Влияние размера группы на качество преимущества]
\label{ex:grpo-group-size}
Предположим, промпт имеет истинное ожидаемое вознаграждение $\mu = 0{,}3$ (модель решает его в 30\% случаев), и вознаграждения бинарные $r_i \in \{0, 1\}$. При $G = 4$ возможные средние по группе равны $\mu_G \in \{0; 0{,}25; 0{,}5; 0{,}75; 1\}$. Вероятность того, что $\mu_G = 0$ (все четыре неверные), равна $(0{,}7)^4 = 0{,}240$, в этом случае все преимущества равны нулю и промпт не вносит никакого обучающего сигнала. При $G = 16$ вероятность того, что все ответы неверны, равна $(0{,}7)^{16} = 0{,}003$\,---\,пренебрежимо мало. Большие группы гарантируют, что почти каждый промпт вносит контрастный сигнал.\qed
\end{example}

% ===========================================================
\section{RLVR: обучение с подкреплением с проверяемыми вознаграждениями}
\label{sec:rlvr}

\index{RLVR!definition}

\subsection{Определение и мотивация}

Многие важные задачи обладают свойством, что правильность ответа можно проверить автоматически, без обученной модели вознаграждения. Математические рассуждения («Равен ли $\frac{1}{3}$ правильному ответу на $\int_0^1 x^2\,dx$?»), генерация кода («Проходит ли эта функция все юнит-тесты?») и формальная логика («Проверяется ли это доказательство по типам?»)\,---\,все они допускают детерминированную верификацию. Этот класс задач мотивирует отдельную парадигму обучения.

\begin{definition}[RLVR]
\label{def:rlvr}
\textbf{Обучение с подкреплением с проверяемыми вознаграждениями} (RLVR) является парадигмой обучения, при которой:
\begin{enumerate}[leftmargin=*]
  \item функция вознаграждения $r(q, o)$ является \emph{детерминированным верификатором}\index{verifier}, проверяющим правильность ответа $o$ на промпт $q$;
  \item верификатор возвращает бинарное вознаграждение: $r(q, o) \in \{0, 1\}$ (или более общее вознаграждение на основе правил);
  \item обученная модель вознаграждения не используется; верификатор\,---\,это фиксированная, доверенная функция.
\end{enumerate}
\end{definition}

RLVR имеет ряд преимуществ перед RLHF с обученной моделью вознаграждения:
\begin{description}[leftmargin=!,labelwidth=3.5cm]
  \item[Нет взлома вознаграждения] Верификатор\,---\,это фиксированная функция истинного ответа, а не обученное приближение. Стратегия не может найти ложные выходные данные, «обманывающие» верификатор (при условии, что верификатор корректен).
  \item[Неограниченные данные] Поскольку верификатор автоматизирован, можно генерировать и оценивать неограниченное число завершений без каких-либо затрат на аннотирование людьми. Это обеспечивает значительно более масштабное онлайн-обучение.
  \item[Чистый сигнал] Бинарные вознаграждения не зашумлены: математический ответ либо верен, либо нет. Нет разногласий между аннотаторами, нет дрейфа калибровки и субъективных суждений.
  \item[Нет обучения модели вознаграждения] Верификатор\,---\,это фиксированная функция (например, символический решатель или тестовая оснастка), поэтому нет модели вознаграждения для обучения, оценки или развёртывания.
\end{description}

\subsection{Формальная постановка}

В рамках RLVR обучающие данные состоят из набора промптов $\mathcal{D} = \{q^{(1)}, \ldots, q^{(N)}\}$, каждый со связанным эталонным ответом $a^{(i)}$. Верификатор $V: (q, o) \to \{0, 1\}$ извлекает ответ модели из завершения $o$ и сравнивает его с $a^{(i)}$:
\begin{equation}
\label{eq:rlvr-reward}
r(q, o) = V(q, o) = \ind[\text{extract}(o) = a].
\end{equation}
Более сложные верификаторы могут присваивать частичный зачёт, проверять промежуточные шаги или объединять несколько сигналов вознаграждения:
\begin{equation}
\label{eq:rlvr-composite}
r(q, o) = w_\mathrm{acc}\, V_\mathrm{accuracy}(q, o) + w_\mathrm{fmt}\, V_\mathrm{format}(o) + w_\mathrm{len}\, V_\mathrm{length}(o),
\end{equation}
где $V_\mathrm{accuracy}$ проверяет правильность, $V_\mathrm{format}$ проверяет соответствие предписанному формату вывода (например, использование тегов \texttt{<think>}\ldots\texttt{</think>}), а $V_\mathrm{length}$ штрафует чрезмерно длинные ответы.

\subsection{Связь с GRPO}

GRPO является естественным алгоритмом оптимизации для RLVR. Бинарная структура вознаграждения гарантирует, что группы содержат смесь правильных и неправильных завершений (для промптов умеренной сложности), производя нетривиальные преимущества. Отсутствие обученной модели вознаграждения означает, что две оставшиеся копии моделей ($\pi_\theta$ и $\pi_\mathrm{ref}$) легко помещаются в видеопамять GPU, а верификатор добавляет пренебрежимо малые вычислительные накладные расходы.

Конвейер обучения DeepSeek-R1 (DeepSeek-AI, 2025) продемонстрировал эту комбинацию в масштабе: GRPO с бинарными вознаграждениями на основе правил, обучаемый непосредственно на базовой языковой модели (без какого-либо этапа дообучения с учителем), породил модель, которая спонтанно развила рассуждение с цепочкой мыслей, самопроверку и стратегии исследования\,---\,способности, возникшие только из обучения RL.

\begin{example}[Эмерджентное рассуждение в DeepSeek-R1-Zero]
\label{ex:r1-zero}
DeepSeek-R1-Zero применил GRPO с бинарными вознаграждениями за точность к базовой языковой модели (без SFT). Во время обучения модель спонтанно выработала несколько сложных поведений:
\begin{itemize}[leftmargin=*]
  \item \textbf{Рассуждение с цепочкой мыслей}: модель начала производить пошаговые выводы, разбивая сложные задачи на подзадачи.
  \item \textbf{Самопроверка}: модель начала проверять свои промежуточные результаты, производя такие фразы, как «Подождите, дайте мне проверить\ldots» перед окончательным ответом.
  \item \textbf{Исследование альтернативных подходов}: когда первый подход не срабатывал, модель начала пробовать альтернативные стратегии решения в рамках одного ответа.
\end{itemize}
Ни одно из этих поведений не вознаграждалось явно\,---\,проверялась только правильность итогового ответа. Они возникли потому, что ответы с цепочкой мыслей с большей вероятностью дают правильные ответы, а сигнал преимущества GRPO подкреплял всю последовательность токенов, приводящую к правильным результатам.\qed
\end{example}

\subsection{Области применения RLVR}

\begin{description}[leftmargin=!,labelwidth=3.5cm]
  \item[Математика] Символическая верификация по сравнению с эталонными ответами. Инструменты такие как \texttt{sympy}, \texttt{math\_verify} и \texttt{latex2sympy} разбирают и сравнивают математические выражения.
  \item[Генерация кода] Наборы юнит-тестов служат верификатором. Модель генерирует код; код выполняется в изолированной среде; вознаграждение\,---\,доля пройденных тестов.
  \item[Формальные доказательства] Средства проверки типов для систем доказательств теорем (Lean, Coq, Isabelle) верифицируют доказательства детерминированно.
  \item[Генерация с ограничениями] Ограничения формата (например, валидация JSON-схемы) предоставляют бинарную обратную связь о структурной правильности.
  \item[Естественные науки] Числовые ответы по физике, химии и инженерии можно верифицировать в пределах допуска.
\end{description}

% ===========================================================
\section{Отбор по отклонению и best-of-$N$}
\label{sec:rejection-sampling}

\index{rejection sampling}\index{best-of-$N$}

До применения GRPO (или любого алгоритма RL) естественным базовым уровнем является сэмплирование \textbf{best-of-$N$}: генерация $N$ завершений, оценка каждого верификатором или моделью вознаграждения и возврат наивысшего по баллам. Иногда это называется \emph{дообучением с отбором по отклонению}, когда лучшие завершения используются как контролируемые обучающие данные.

\subsection{Алгоритм}

Для каждого промпта $q$ сэмплируем $o_1, \ldots, o_N \sim \pi_\theta(\cdot \mid q)$ и выбираем $o^* = \argmax_{i} r(q, o_i)$. Ожидаемое вознаграждение best-of-$N$ растёт с $N$: если каждое завершение правильно с вероятностью $p$ (бинарное вознаграждение), вероятность того, что хотя бы одно из $N$ правильно, равна $1 - (1-p)^N$.

\begin{example}[Best-of-$N$ против одного сэмпла]
\label{ex:best-of-n}
Модель решает данную математическую задачу с вероятностью $p = 0{,}3$. С одним сэмплом ожидаемое вознаграждение равно 0,3. При $N = 8$ вероятность хотя бы одного правильного решения равна $1 - (0{,}7)^8 = 0{,}942$. При $N = 64$ она возрастает до $1 - (0{,}7)^{64} > 0{,}9999$. Best-of-$N$ драматически улучшает производительность во время инференса, но стоимость линейно растёт с $N$.\qed
\end{example}

\subsection{Дообучение с отбором по отклонению}

Вместо использования best-of-$N$ только во время инференса можно использовать отобранные завершения как обучающие данные SFT. Этот процесс, называемый дообучением с отбором по отклонению (RFT), создаёт улучшенный набор данных из собственных генераций модели:
\begin{enumerate}[leftmargin=*]
  \item Генерировать $N$ завершений на промпт.
  \item Фильтровать, оставляя только правильные (или высокооцениваемые) завершения.
  \item Дообучить $\pi_\theta$ на отфильтрованном наборе данных с использованием стандартного обучения с учителем.
\end{enumerate}

\subsection{Сравнение с GRPO}

Best-of-$N$ и GRPO тесно связаны. Оба сэмплируют несколько завершений и используют вознаграждение для различения хороших и плохих. Ключевые различия:
\begin{itemize}[leftmargin=*]
  \item \textbf{Использование сигнала.} Best-of-$N$ отбрасывает всю информацию, кроме единственного лучшего завершения. GRPO использует полное распределение вознаграждений для вычисления преимуществ, извлекая контрастный обучающий сигнал из каждого завершения.
  \item \textbf{Качество градиента.} Best-of-$N$ сводится к SFT на отобранных примерах, что не может снизить вероятность плохих ответов. GRPO явно снижает вероятность отрицательно-адвантажных завершений.
  \item \textbf{Итеративное улучшение.} Best-of-$N$ с RFT требует чередования этапов генерации и SFT. GRPO интегрирует генерацию и оптимизацию стратегии в единый цикл.
  \item \textbf{Когда отбора по отклонению достаточно.} Если модель уже решает задачу с умеренной вероятностью ($p \geq 0{,}3$) и цель\,---\,качество во время развёртывания, best-of-$N$ при инференсе может быть более экономически эффективным, чем обучение RL. Если модель редко решает задачу ($p < 0{,}05$) и цель\,---\,улучшить базовую модель, необходимо обучение RL (GRPO) для сдвига распределения.
\end{itemize}

% ===========================================================
\section{Онлайн-обучение RL в масштабе}
\label{sec:grpo-scale}

\index{GRPO!training at scale}

\subsection{Двухфазная архитектура}

Производственное обучение GRPO в масштабе разделяет нагрузку на две асинхронные фазы:
\begin{description}[leftmargin=!,labelwidth=3.5cm]
  \item[Фаза генерации] Работники роллаута размещают копии текущей стратегии $\pi_\theta$ для инференса и генерируют завершения. Эта фаза ограничена памятью (веса модели) и чувствительна к задержке (авторегрессивная генерация токенов). Для максимизации пропускной способности используются оптимизированные движки инференса, такие как vLLM (Kwon и др., 2023), с непрерывным батчингом и PagedAttention.
  \item[Фаза обучения] Работники обучения получают сгенерированные завершения, вычисляют преимущества и выполняют обновления градиента. Эта фаза ограничена вычислениями (прямые и обратные проходы через полную модель с вычислением градиента).
\end{description}
Две фазы могут работать в конвейерном режиме: пока обучающий модуль обновляется на батче $k$, работники роллаута генерируют батч $k+1$ с использованием последней контрольной точки. Это перекрытие скрывает задержку генерации и поддерживает высокую утилизацию GPU.

\subsection{Обработка завершений переменной длины}

Практической проблемой является то, что завершения внутри группы могут резко различаться по длине\,---\,от нескольких токенов («1/3») до тысяч токенов (многостраничный вывод с цепочкой мыслей). Используются несколько стратегий:
\begin{itemize}[leftmargin=*]
  \item \textbf{Дополнение и маскировка.} Дополнить все завершения в группе до максимальной длины и использовать маски внимания. Просто, но тратит память на короткие завершения.
  \item \textbf{Упаковка.} Конкатенировать завершения из разных промптов в единую последовательность, разделённую токенами-разделителями, и использовать блочно-диагональные маски внимания для предотвращения перекрёстного смешения. Это максимизирует утилизацию GPU.
  \item \textbf{Максимальная длина генерации.} Применять жёсткое ограничение на длину завершения. Для задач рассуждения обычны ограничения в 4096\,---\,8192 токена.
\end{itemize}
Нормировка $1/|o_i|$ в цели GRPO~\eqref{eq:grpo-objective} необходима при завершениях переменной длины: без неё длинные завершения вносили бы непропорционально большой вклад в потерю.

\subsection{Формирование батча}

Для GRPO каждый «мини-батч» состоит из промптов, каждый с $G$ завершениями. Если батч содержит $B$ промптов, общее число завершений равно $BG$. DeepSeek-R1 использует $BG = 8192$ завершения на роллаут, которые случайным образом разбиваются на 16 мини-батчей по 512 завершений для одной эпохи обучения.

Стоимость генерации на итерацию масштабируется как $O(BG \cdot \bar{T})$, где $\bar{T}$\,---\,средняя длина завершения. Для задач долгосрочного рассуждения ($\bar{T} \approx 2000$ токенов) при $BG = 8192$ это примерно 16 миллионов токенов на роллаут\,---\,существенная нагрузка инференса, которая мотивирует использование выделенных кластеров инференса.

\subsection{Синхронный и асинхронный дизайн}

В \textbf{синхронном} дизайне обучающий модуль ждёт завершения генерации всеми работниками роллаута, прежде чем начать фазу обучения. Это гарантирует, что все завершения из одной и той же контрольной точки стратегии, делая отношения важности $\rho_{i,t} = 1$ в начале каждой фазы обучения. Недостатком является простой: обучающие модули ждут во время генерации, а работники роллаута\,---\,во время обучения.

В \textbf{асинхронном} дизайне работники роллаута непрерывно генерируют завершения и помещают их в очередь, тогда как обучающие модули извлекают из очереди и обновляют стратегию. Стратегия, использованная для генерации, может отставать от обучаемой стратегии на один или несколько шагов. Это вносит лёгкое несоответствие вне текущей стратегии, но отсечённое отношение важности в GRPO естественно с этим справляется: если $\rho_{i,t}$ слишком сильно отклоняется от 1, отсечение ограничивает вклад в градиент, предотвращая нестабильность из-за устаревших данных. Асинхронные дизайны достигают значительно более высокой утилизации GPU и предпочтительны в масштабе.

\subsection{Воспроизведение опыта для GRPO}

\index{experience replay!GRPO}
В отличие от PPO, который обычно использует данные по текущей стратегии несколько эпох, а затем отбрасывает их, GRPO может извлечь пользу из лёгкой формы воспроизведения опыта. Недавние завершения из предыдущих $M$ итераций можно хранить в буфере воспроизведения и смешивать со свежими данными по текущей стратегии. Отношение важности $\rho_{i,t}(\theta)$ естественно корректирует устаревание воспроизводимых данных, а механизм отсечения предотвращает чрезмерные коррекции вне текущей стратегии. Воспроизведение снижает стоимость генерации, амортизируя каждое завершение на несколько шагов обучения, ценой незначительно возросшего несоответствия распределений.

\subsection{Мониторинг и диагностика}

Эффективное обучение GRPO требует мониторинга нескольких величин:
\begin{description}[leftmargin=!,labelwidth=3.5cm]
  \item[Кривая вознаграждения] Среднее вознаграждение по промптам должно расти в ходе обучения. Плато может указывать на то, что распределение промптов слишком лёгкое (все промпты уже решены) или слишком сложное (нет контрастного сигнала).
  \item[Дивергенция KL] Скользящее среднее $\KL{\pi_\theta}{\pi_\mathrm{ref}}$ должно расти медленно. Быстрый рост KL сигнализирует об опасном удалении стратегии от референса; следует увеличить $\beta$ или уменьшить скорость обучения.
  \item[Доля отсечения] Доля токенов, где отношение важности $\rho_{i,t}$ отсечено. Доля отсечения выше 10\,---\,15\% указывает на то, что стратегия меняется слишком быстро за итерацию; следует рассмотреть уменьшение скорости обучения или числа внутренних эпох $K_e$.
  \item[Разнообразие групп] Доля групп, содержащих как положительные, так и отрицательные преимущества. Если эта доля падает ниже 20\%, промпты могут быть слишком лёгкими или слишком сложными, и следует перебалансировать распределение промптов.
  \item[Величина преимущества] Средняя абсолютная величина преимущества $|\hat{A}_i|$ должна быть приблизительно равна 1 (по построению, благодаря стандартизации). Экстремальные значения указывают на численные проблемы при нормализации.
\end{description}

% ===========================================================
\section{Сравнение: PPO, GRPO и DPO}
\label{sec:grpo-comparison}

Таблица~\ref{tab:ppo-grpo-dpo} содержит детальное сравнение трёх основных алгоритмов выравнивания и RL, рассмотренных в главах~\ref{ch:actor-critic}, \ref{ch:dpo} и настоящей главе.

\begin{table}[ht]
\centering
\caption{Детальное сравнение PPO, GRPO и DPO для обучения LLM.}
\label{tab:ppo-grpo-dpo}
\footnotesize
\begin{tabular}{@{}lp{3.2cm}p{3.2cm}p{3.2cm}@{}}
\toprule
\textbf{Свойство} & \textbf{PPO} (гл.~\ref{ch:actor-critic}) & \textbf{GRPO} (эта глава) & \textbf{DPO} (гл.~\ref{ch:dpo}) \\
\midrule
Онлайн / автономный & Онлайн & Онлайн & Автономный \\
Нужен критик? & Да (сеть функции ценности) & Нет & Нет \\
Модель вознаграждения? & Да (обученная RM) & Верификатор или RM & Нет (неявная в потере) \\
Копии моделей & 4 (актор, критик, реф., RM) & 2 (актор, реф.) & 2 (актор, реф.) \\
Оценка преимущества & GAE (на токен) & Групповая нормализация (на последовательность) & Нет (без преимущества) \\
Доверительная область & Отсечённое отношение & Отсечённое отношение & Неявная KL \\
Требование к данным & Только промпты & Только промпты & Пары предпочтений \\
Исследование & Сэмплирование по текущей стратегии & Сэмплирование по текущей стратегии & Нет (фиксированный набор данных) \\
Сдвиг распределения & Нет & Нет & Да (деградирует со временем) \\
Присваивание заслуг & На токен & На последовательность & На пару \\
Риск взлома вознаграждения & Высокий (обученная RM) & Низкий (верификатор) & Нет (нет RM) \\
Сложность реализации & Высокая & Средняя & Низкая \\
\bottomrule
\end{tabular}
\end{table}

\subsection{Когда применять каждый метод}

\textbf{Применяйте PPO}, когда сигнал вознаграждения исходит от обученной модели вознаграждения и важно присваивание заслуг на уровне токена\,---\,например, в общем RLHF, где качество отдельных токенов влияет на общий ответ. Основанные на GAE преимущества PPO обеспечивают тонкий сигнал, но инфраструктурные затраты высоки.

\textbf{Применяйте GRPO}, когда вознаграждения проверяемы или когда приоритетны простота и экономия памяти. GRPO является предпочтительным алгоритмом для задач рассуждения (математика, код, логика), где доступны бинарные вознаграждения за правильность. Он также эффективен с обученными моделями вознаграждения, когда стоимость критика чрезмерна.

\textbf{Применяйте DPO}, когда имеются данные предпочтений, но онлайн-генерация слишком дорога или непрактична. DPO является самым быстрым путём от SFT-модели к выровненной модели, не требуя инфраструктуры RL. Для задач, где сдвиг распределения вызывает озабоченность, итеративный DPO (сбор новых данных предпочтений между раундами) частично закрывает разрыв.

\begin{remark}
Эти методы не являются взаимоисключающими. Общий производственный паттерн\,---\,начальный запуск с DPO на имеющемся наборе данных предпочтений, затем уточнение с GRPO на проверяемых задачах и, наконец, шлифовка с PPO с использованием обученной модели вознаграждения для открытого качества. Каждый этап устраняет отдельное измерение выравнивания (см. пример~\ref{ex:alignment-choice} в главе~\ref{ch:dpo}).
\end{remark}

\begin{example}[Бюджет памяти при масштабе 70B]
\label{ex:grpo-memory}
Рассмотрим обучение модели с 70B параметрами в формате половинной точности (2 байта на параметр, то есть 140 ГБ на копию модели). PPO-RLHF требует четыре копии: актор (140 ГБ), критик (140 ГБ), референс (140 ГБ) и модель вознаграждения (140 ГБ), итого 560 ГБ без учёта состояний оптимизатора, активаций или градиентов. Один узел 8$\times$A100 (640 ГБ всего) едва вмещает модели, оставляя почти не места для активаций.

GRPO-RLVR требует две копии: актор (140 ГБ) и референс (140 ГБ), итого 280 ГБ. Верификатор (например, символический математический проверщик) пренебрежимо мал. Это оставляет 360 ГБ на том же узле для состояний оптимизатора и активаций\,---\,достаточно для комфортного обучения с контрольными точками градиента. Двукратное снижение требований к памяти нередко является решающим фактором при выборе GRPO вместо PPO в большом масштабе.\qed
\end{example}

% ===========================================================
\section{Итоги}
\label{sec:grpo-summary}

В этой главе была разработана групповая относительная оптимизация стратегии (GRPO)\,---\,алгоритм градиента стратегии, устраняющий критика с функцией ценности путём оценки преимуществ по группам сэмплированных завершений.

Преимущество GRPO (определение~\ref{def:grpo-advantage}) стандартизирует вознаграждение каждого завершения относительно среднего и стандартного отклонения по группе, производя сигнал с нулевым средним и приблизительно единичной дисперсией, не требующий обученных параметров. Отсечённая суррогатная цель (уравнение~\eqref{eq:grpo-clip-token}) наследует свойство доверительной области PPO, обеспечивая безопасное обучение с несколькими эпохами. Штраф KL (уравнение~\eqref{eq:grpo-kl}) применяется непосредственно в функции потерь, а не во вознаграждении, чётко разделяя регуляризацию и вычисление преимущества.

Теоретический анализ (раздел~\ref{sec:grpo-theory}) показал, что градиент GRPO тесно связан с REINFORCE с базисом среднего по группе, со смещением $O(1/G)$, исчезающим при умеренных размерах группы. Стандартизация на $\sigma_G$ автоматически нормирует величины градиентов по промптам с разной дисперсией вознаграждений\,---\,практическое преимущество перед RLOO (раздел~\ref{sec:rloo} главы~\ref{ch:pg}).

RLVR (раздел~\ref{sec:rlvr}) формализовал парадигму обучения с детерминированными верификаторами вместо обученных моделей вознаграждения. Комбинация GRPO и RLVR обеспечивает крупномасштабный онлайн-RL с чистыми сигналами вознаграждения, без взлома вознаграждения и с минимальными инфраструктурными накладными расходами\,---\,конфигурация, породившая DeepSeek-R1 и аналогичные модели рассуждений.

Таблица сравнения (таблица~\ref{tab:ppo-grpo-dpo}) поместила GRPO в контекст: он занимает онлайн, безкритиковое пространство проектирования, дополняя PPO (онлайн, с критиком) и DPO (автономный, без критика).

\begin{description}[leftmargin=!,labelwidth=3.5cm]
  \item[Глава~\ref{ch:pg}] закладывает фундаментальный оценщик REINFORCE и базис RLOO, на которых строится GRPO.
  \item[Глава~\ref{ch:actor-critic}] развивает PPO\,---\,алгоритм, отсечённый суррогат которого GRPO наследует, но критика которого устраняет.
  \item[Глава~\ref{ch:dpo}] развивает DPO\,---\,автономную альтернативу, проблема сдвига распределения которой мотивирует онлайн-подход GRPO.
  \item[Глава~\ref{ch:practical-rlhf}] описывает полный производственный конвейер RLHF, где GRPO может заменить PPO в качестве оптимизатора RL.
\end{description}

% ===========================================================
\section*{Упражнения}
\addcontentsline{toc}{section}{Упражнения}

\begin{exercise}
\label{ex:grpo-1}
\textbf{Преимущество GRPO вручную.} Промпт получает $G = 5$ завершений с вознаграждениями $r = (0, 1, 1, 0, 1)$. Вычислите групповое среднее $\mu_G$, групповое стандартное отклонение $\sigma_G$ и преимущество GRPO $\hat{A}_i$ для каждого завершения. Проверьте, что $\sum_i \hat{A}_i \neq 0$, но $\sum_i \hat{A}_i \approx 0$ (объясните, почему точный нуль не гарантирован при использовании среднего по полной группе).
\end{exercise}

\begin{exercise}
\label{ex:grpo-2}
\textbf{GRPO против RLOO.} Используя вознаграждения из упражнения~\ref{ex:grpo-1}, вычислите преимущество RLOO $\hat{A}_i^{\mathrm{RLOO}} = \frac{K}{K-1}(r_i - \bar{r})$ для каждого завершения. Сравните два набора преимуществ. При каком условии на распределение вознаграждений два оценщика дали бы одно и то же \emph{направление} градиента?
\end{exercise}

\begin{exercise}
\label{ex:grpo-3}
\textbf{Вырожденные группы.} Что происходит с преимуществом GRPO, когда все завершения в группе получают одинаковое вознаграждение? Почему такое поведение желательно с точки зрения обучения? Предложите модификацию, которая всё же извлекала бы сигнал из таких групп.
\end{exercise}

\begin{exercise}
\label{ex:grpo-4}
\textbf{Смещение базиса среднего по группе.} Докажите, что смещение от использования $\mu_G = \frac{1}{G}\sum_j r_j$ в качестве базиса (вместо среднего leave-one-out) в оценщике REINFORCE точно равно $\frac{1}{G}\nabla_\theta J(\theta)$. Сделайте вывод, что смещение исчезает при $G \to \infty$.
\end{exercise}

\begin{exercise}
\label{ex:grpo-5}
\textbf{Вероятность best-of-$N$.} Модель имеет точность $p$ на единичный сэмпл для данной задачи. Выведите ожидаемое вознаграждение (точность) сэмплирования best-of-$N$ как функцию от $p$ и $N$. Для $p = 0{,}1$ вычислите значение $N$, необходимое для достижения 90\% ожидаемой точности.
\end{exercise}

\begin{exercise}
\label{ex:grpo-6}
\textbf{Свойства оценщика KL.} Покажите, что оценщик KL GRPO~\eqref{eq:grpo-kl} неотрицателен и равен нулю тогда и только тогда, когда $\pi_\theta = \pi_\mathrm{ref}$ в данной позиции токена. \emph{Подсказка:} используйте неравенство $e^x \geq 1 + x$.
\end{exercise}

\begin{exercise}
\label{ex:grpo-7}
\textbf{Размер группы и покрытие.} Для бинарных вознаграждений с точностью $p$ на единичный сэмпл выведите вероятность того, что группа размера $G$ содержит \emph{как} правильные, так и неправильные завершения (то есть производит нетривиальный контрастный сигнал). Постройте эту вероятность как функцию от $G$ для $p \in \{0{,}1; 0{,}3; 0{,}5; 0{,}7; 0{,}9\}$.
\end{exercise}

\begin{exercise}
\label{ex:grpo-8}
\textbf{GRPO с вознаграждением за процесс.} Стандартный GRPO присваивает одинаковое преимущество каждому токену в завершении (надзор за результатом). Предположим, верификатор процесса присваивает вознаграждение $r_i^{(t)}$ после каждого шага рассуждения $t$. Предложите модификацию преимущества GRPO, включающую вознаграждения на уровне шагов, сохраняя при этом отсутствие обученного критика. Запишите модифицированную цель.
\end{exercise}

\begin{exercise}
\label{ex:grpo-9}
\textbf{Эффективность отбора по отклонению против GRPO.} Задача имеет точность $p = 0{,}2$ на единичный сэмпл. Сравните следующие две стратегии улучшения модели: (a)~дообучение с отбором по отклонению best-of-$N$ при $N = 16$, обучение на правильных завершениях одну эпоху; (b)~GRPO с $G = 16$. Какая стратегия использует сгенерированные данные более эффективно? Обсудите компромиссы с точки зрения использования данных, качества градиента и стоимости вычислений.
\end{exercise}

\begin{exercise}
\label{ex:grpo-10}
\textbf{Реализация.} Реализуйте упрощённый цикл обучения GRPO в виде псевдокода на PyTorch. Ваша реализация должна: (i)~сэмплировать $G$ завершений на промпт, (ii)~вычислять групп-нормированные преимущества, (iii)~вычислять отсечённую суррогатную потерю, и (iv)~добавлять штраф KL. Предполагается наличие функций \texttt{model.generate()}, \texttt{model.log\_probs()}, \texttt{ref\_model.log\_probs()} и \texttt{reward\_fn()}.
\end{exercise}
